\clearpage
\section{Planificación del proyecto}
El programa combina cuatro frentes paralelos: parche descartable, lector reutilizable, pulsera y aplicación. Compras, integración, QA, clínica, regulación y captación de aliados son actividades transversales. El desarrollo dura doce meses y el cierre regulatorio se prevé en el mes 18. El prototipo integrado del mes 7 habilita la revisión de aptitud y el avance de presentaciones según la evidencia disponible.
\subsection{Equipo, capacidad y sprints}
Se mantienen dos perfiles de ingeniería y dos técnicos/administrativos hasta el mes 18, a 160 horas mensuales por persona: 7.680 horas el primer año y 3.840 horas adicionales. Las horas controlan capacidad y no se vuelven a sumar como salarios por tarea. Los cuatro recursos se comparten entre frentes y realizan revisión cruzada. Laboratorios e investigadores clínicos se contratan como servicios externos.
Se planifican 26 sprints de dos semanas el primer año y 13 durante los meses 13--18. La asignación mensual de capacidad se ajusta a los dos o tres sprints de cada mes; no se multiplica automáticamente cada sprint por 80 horas. Se asigna el 80\% de capacidad a los cuatro productos y el 20\% a coordinación, compras, QA, clínica, regulación y captación de aliados. Las 195 tareas de los 39 sprints contienen estimaciones preliminares ajustables en cada revisión. Cada tarea incluye ID, producto, dependencia, sprint, responsable, horas por recurso, entregable y aceptación en la hoja Tareas y sprints. La compra, recepción, revisión externa y autorización conservan sus esperas independientes del esfuerzo.
\begin{center}\begin{tabular}{lll}\hline
Meses & Sprints & Resultado\\\hline
1--3 & 1--6 & Requisitos, interfaces, proveedores\\
4--6 & 7--13 & Subsistemas construidos y ensayados\\
7 & 14--15 & Prototipo integrado\\
8--10 & 16--22 & Pruebas, correcciones y avance regulatorio\\
11--12 & 23--26 & MVP técnico y documentación\\
13--18 & 27--39 & Evaluación, registro y lanzamiento preparado\\\hline
\end{tabular}\end{center}
\subsection{Secuencia por producto}
Parche: materiales y reactivos, muestras, calibración e interferencias, estabilidad y adhesión, fabricación y control de lotes. Lector: interfaces, componentes y PCB, montaje, firmware y adquisición, memoria y BLE, USB-C y carcasa, pruebas. Pulsera: sensores y PCB, montaje, firmware, BLE directo y LED, alimentación y protección, pruebas. Aplicación: interfaces y datos, BLE, historial, calidad y recuperación, avisos actuales, predicciones y ayuda. La integración exige versiones compatibles y resultados individuales disponibles. Cada sprint cierra con demostración, revisión de riesgos y evidencia.
\subsection{Evaluación clínica y autorizaciones}
La preparación de requisitos, asesoramiento y documentación comienza en los meses 1--6. Desde el prototipo del mes 7 se avanza con los ensayos y presentaciones permitidos por la evidencia de seguridad. Participan ANMAT, el comité de ética y el centro clínico para autorizar la investigación; se distingue esa autorización del registro para comercializar. Se incluye el encuadre de la aplicación, fabricación/importación y revisión de homologación BLE ante ENACOM. Los estudios con participantes se realizan después de las autorizaciones correspondientes. Se prevén evaluación de mediciones y alertas frente a referencia, resultados, observaciones, correcciones y documentación; el presente plan no desarrolla un protocolo clínico completo.
El cierre regulatorio hacia el mes 18 es una previsión ajustable. La captación de médicos y clínicas avanza desde el mes 9 con responsables técnicos/administrativos y asesoramiento externo. Las ventas comienzan desde el mes 19, sujetas a autorizaciones. Las unidades de investigación se distinguen de los 100 kits comerciales previstos para el lanzamiento.
\subsection{Control y aceptación}
La revisión mensual concilia tareas, capacidad, compras, costos y riesgos. Se comprueban adquisición, calibración, interferencias, autonomía, memoria y retención, BLE simultáneo, recuperación, recambio de lector, carga, temperatura, IP68 y sudor, LED con datos válidos o vencidos y trazabilidad de lotes. Los meses 3, 6, 7, 10, 12 y 18 tienen revisión formal. Los cambios afectan simultáneamente backlog, cronograma, presupuesto, documentación y evidencia clínica cuando corresponda.
